\setauthor{\secondauthor}
\section{Analyse der Chatbot Nutzerdaten}

Zur Evaluation des Chatbots standen sowohl quantitative als auch qualitative 
Daten zur Verfügung. Die quantitativen Daten ergaben sich aus den protokollierten 
Sitzungen mit dem Chatbot. Für die untersuchung wurden 30 individuelle Sitzungen herangezogen.
Hierzu zählen insbesondere die Anzahl der 
Konversationen, die Verteilung der getesteten Varianten, die Sitzungsdauer, die 
Anzahl der Nachrichten, das geäußerte Download-Interesse sowie die abgegebenen 
Feedback-Bewertungen. Die Transcripts aus den Konversationen wurden qualitativ 
ausgewertet, um Rückschlüsse auf häufige Fragen, Unsicherheiten und immer wieder 
geäußerte Bedenken von potentiellen Nutzer:innen zu ziehen.

\subsection{Einfluss der Personalisierung auf die wahrgenommene Relevanz}

Zur Untersuchung der Rolle von Personalisierung wurden die Chatbot-Varianten
mit stärker personalisierter Ansprache den Varianten mit geringerer
Personalisierung gegenübergestellt. Personalisierung wurde im Chatbot
insbesondere durch die Verwendung des Namens, die Unterscheidung zwischen
Zielgruppen sowie die angepasste Ansprache umgesetzt.

Die Auswertung zeigt, dass personalisierte Varianten im Durchschnitt
\textbf{3,3/5}
Punkte erreichten, während die weniger personalisierten Varianten
durchschnittlich \textbf{2,9/5}
Punkte erzielten. Auch im Hinblick auf das Download-Interesse ergab sich ein
Unterschied: In \textbf{35\%}
der personalisierten Sitzungen wurde ein konkretes Download-Interesse
festgestellt, während dies bei den weniger personalisierten Varianten in
\textbf{25\%} der Fälle zutraf.

Diese Ergebnisse deuten darauf hin, dass Personalisierung einen positiven
Einfluss auf die wahrgenommene Relevanz des Chatbots haben kann. Besonders
deutlich wird dies, wenn Nutzer:innen in ihren Rückmeldungen Bezug auf die
für sie passende Ansprache oder auf die zielgruppenspezifische Erklärung der
App nahmen. Beispiele dafür finden sich in den Transcripts, etwa wenn
Nutzer:innen Formulierungen verwendeten wie
\textbf{"Ich bin Jobsuchender was bringt mir die App?"}.

Gleichzeitig zeigte sich, dass Personalisierung nicht in jedem Fall zu
deutlich besseren Ergebnissen führte. In einzelnen Sitzungen wirkte die
Ansprache eher neutral oder wurde von den Nutzer:innen nicht explizit
aufgegriffen. Daraus lässt sich ableiten, dass Personalisierung zwar ein
relevanter Faktor ist, ihre Wirkung jedoch vom Gesamtkontext der
Interaktion abhängt.

\subsection{Rolle von Design und Sprache für das Nutzererlebnis}

Neben der Personalisierung wurden auch die Gestaltung und Sprache des 
Chatbots untersucht. Denn neben der Qualität der Inhalte, für die der Bot steht, 
wird das Nutzererlebnis stark auch von der Form und Strukturierung der Inhalte 
sowie von der Verständlichkeit des Dialogs beeinflusst. Die Varianten 
unterschieden sich vor allem in der sprachlichen Tonalität, während die zugrunde 
liegende Konversationslogik meist identisch war. Anhand der Transcripts und 
Nutzungsdaten wurde dann auch untersucht, ob bestimmte Formulierungen, 
Dialogpfade oder Antwortmuster mit längeren Gesprächen oder besserem Feedback 
korrelierten.

Die durchschnittliche Gesprächsdauer betrug bei den einzelnen Varianten:
\begin{itemize}
    \item Variante 1: \textbf{2 min 18 s}
    \item Variante 2: \textbf{1 min 49 s}
    \item Variante 3: \textbf{2 min 05 s}
    \item Variante 4: \textbf{1 min 37 s}
\end{itemize}

Die durchschnittliche Anzahl an Nachrichten inklusive Nachrichten des Agent pro Sitzung lag bei:
\begin{itemize}
    \item Variante 1: \textbf{5,6}
    \item Variante 2: \textbf{4,7}
    \item Variante 3: \textbf{5,2}
    \item Variante 4: \textbf{4,5}
\end{itemize}

Diese Kennzahlen können nun als nächster Schritt dazu dienen, zu beurteilen, wie 
sehr Nutzer:innen sich auf die Konversation eingelassen haben. Lange oder viele 
Sitzungen sind also nicht zwangsläufig gut, können aber auf höheres Interesse 
oder mehr Informationsbedarf hindeuten. Besonders interessant werden diese Werte 
im Vergleich zu Feedback und Gesprächsinhalten.

Qualitativ hat sich in den Transcripts gezeigt, dass positive Nutzererlebnisse 
besonders dann zu verzeichnen waren, wenn der Chatbot klare, kurze, verständliche 
Antworten gab. Negativer wurde es dann, wenn Antworten zu allgemein, zu ungenau 
oder zu unpersönlich waren.

Insgesamt kann man feststellen, dass informelle Varianten zwar leicht positiver
bewertet wurden, formelle Varianten jedoch häufiger zu konkretem Download-Interesse 
geführt haben. Somit lässt sich keine eindeutige Überlegenheit eines bestimmten
Stils ausmachen.

\subsection{Vergleich von freundlicher und formeller Tonalität}

Ein zentrales Untersuchungsanliegen war die Frage, ob ein freundlicher,
lockerer Ton von den Nutzer:innen besser angenommen wird als ein formeller
Sprachstil. Dazu wurden die informellen Varianten den formelleren Varianten
gegenübergestellt.

Die durchschnittliche Zufriedenheit bei den informellen Varianten liegt bei
\textbf{3,2/5}, während die formellen Varianten einen
Durchschnittswert von \textbf{3,0/5} erreichten. Auch bei der
Download-Intention ergaben sich Unterschiede: In
\textbf{23\%} der informellen Sitzungen zeigte
sich Download-Interesse, bei den formellen Varianten in
\textbf{37\%}.

Zusätzlich wurden die freien Rückmeldungen betrachtet. Dabei zeigte sich,
dass der informelle Ton in mehreren Fällen als \textbf{"weniger kompetent" oder "freundlich"}
wahrgenommen wurde. Der formelle Stil wurde dagegen eher mit
\textbf{"umständliches FAQ" und "zu komplex"}
assoziiert.

\subsection{Häufige Fragen und Bedenken vor dem Download}

Um feststellen zu können welche Bedenken häufig vor einem Download geäußert werden,
wurden die Gesprächs-Transskripte analysiert und häufige oder ähnliche Fragen 
zusammengefasst.

Die häufigsten Fragen bezogen sich auf:
\begin{itemize}
    \item \textbf{Funktionsweise der App}
    \item \textbf{Einsatzbereich}
    \item \textbf{Registrierung}
\end{itemize}

Die qualitativen Auswertungen zeigen damit, dass ein großer Teil der 
Konversationen nicht nur allgemeine Interessensbekundungen beinhalteten, sondern 
auch konkrete Informationsbedarfe und Unsicherheiten sichtbar machten. Die 
Ergebnisse sind für die Optimierung des Chatbots besonders wertvoll, da sie 
zeigen, welche Inhalte künftig noch klarer oder früher im Gespräch adressiert 
werden sollten.

\subsection{Bewertung der Tracking-Werkzeuge}

Neben der direkten Nutzerinteraktion wurde auch untersucht, welche Werkzeuge
sich gut zur Analyse des Nutzerverhaltens eignen. Für das Projekt standen
dabei insbesondere Voiceflow-interne Werte, Microsoft Clarity und Google
Analytics zur Verfügung.
Voiceflow lieferte vor allem Einblick in die eigentliche
Konversation, etwa in Gesprächsverläufe, Variantenzuordnung und
konkrete Nutzereingaben. Diese Werte waren besonders wertvoll für die
qualitative Analyse der Chatbot-Interaktion sowie für den direkten
Vergleich der verschiedenen Chatbot-Varianten.
Microsoft Clarity erlaubte darüber hinaus mehr Einblick in das konkrete
Verhalten auf der Website, etwa durch Session Recordings oder Heatmaps, und
damit auch in das Verhalten der Nutzer:innen vor, in und nach der Nutzung
des Chatbots auf der Seite. Google Analytics lieferte
strukturiert Kennzahlen zu Aufrufen, Sitzungen und allgemeinen
Nutzungspfaden.
Zur Evaluierung des Projekts stellte sich die Kombination aus
Voiceflow-Daten, Microsoft Clarity und Google Analytics als am
hilfreichsten heraus. Voiceflow war für die eigentliche Chatbot-Analyse
unersetzbar, da dort die inhaltlichen und konversationsbezogenen Daten
vorliegen. In der vergleichenden Analyse des Nutzerverhaltens auf der gesamten
Website erwiesen sich jedoch Microsoft Clarity und Google Analytics als
hilfreicher, da sie mehr Vergleichbarkeit zwischen Chatbot und
Multi-Step-Funnel sowie präzisere Aufzeichnungen des realen
Nutzerverhaltens ermöglichen.

\subsection{Zwischenfazit}

Die Auswertung der Chatbot-Nutzerdaten zeigt, dass der Chatbot prinzipiell 
geeignet ist, Nutzer:innen strukturiert über die App \textit{Jobbing} zu informieren und 
dabei nützliche Daten zur Evaluation zu liefern. Besonders deutlich zeigte sich, 
dass Personalisierung die wahrgenommene Relevanz des Systems positiv beeinflussen 
kann. Die personalisierten Varianten erzielten sowohl bessere Zufriedenheitswerte 
als auch ein höheres Downloadinteresse, was darauf hindeutet, dass eine stärker 
zielgruppenorientierte Ansprache im Marketingkontext sinnvoll ist.

In Bezug auf die Tonalität ergab sich kein eindeutiges Bild. Während die 
informell gehaltenen Varianten im Durchschnitt etwas besser bewertet wurden, 
zeigte sich bei den formell gehaltenen Varianten ein höheres konkretes 
Downloadinteresse. Daraus folgt, dass die sprachliche Ausgestaltung nicht 
isoliert betrachtet werden darf, sondern immer in Verbindung mit 
Personalisierung, Inhalt und Gesprächsverlauf bewertet werden muss. Für die 
Ausgestaltung des Chatbots bedeutet dies also, dass weder der rein lockere noch 
der ausschließlich formelle Stil automatisch überlegen ist. Vielmehr ist eine 
ausgewogene und zielgerichtete Kommunikation entscheidend.

Ebenfalls zeigt sich hier, dass die Kombination unterschiedlicher 
Trackingwerkzeuge nötig ist, um sowohl die dialogische Nutzung des Chatbots als 
auch das allgemeine Nutzungsverhalten umfassend bewerten zu können. So stellt die 
Analyse des Chatbots also eine wichtige Grundlage für den späteren Vergleich mit 
dem Multi-Step-Funnel dar.

\section{Funnel}

\section{Vergleich}
